% LINA · Relatório institucional (capa, página interna, finalização).
% Compile com XeLaTeX. Só comandos de lina.cls: nada de \vspace, \fontsize, \color, \hspace ou \newpage aqui.
\documentclass[relatorio]{lina}

\referencia{REL-2026-01}
\versao{1.0}
\data{10/2026}
\titulocurto{Relatório anual de atividades}
\siglarodape{LINA · UFC}

\begin{document}

% ---------- Capa ----------
\begin{capa}{LINA · UFC\newline Fortaleza, 2026}
  \rotulo{Relatório institucional}
  \titulo{Relatório anual de atividades 2026}
  \subtitulo{Pesquisa, ensino e extensão do Laboratório de Modelagem da Informação}
  \begin{metadados}[18]{YYR}
    \destaque{Laboratório de Modelagem da Informação}\newline Arte · Arquitetura · Ambiente &
    Departamento de Arquitetura, Urbanismo e Design\newline Universidade Federal do Ceará &
    \mostrar{referencia}\newline v. \mostrar{versao}
  \end{metadados}
\end{capa}

% ---------- Página interna ----------
\section{1 · Apresentação}
\begin{multicols}{2}
O Laboratório de Modelagem da Informação reúne pesquisadores de arquitetura, urbanismo e design em torno da informação que descreve o espaço, da escala do objeto à da paisagem. Este relatório apresenta as atividades do ano nas cinco linhas de pesquisa do grupo.

Cada seção descreve projetos em andamento, produção bibliográfica, orientações e parcerias. Os indicadores ao final seguem os critérios da plataforma institucional e permitem comparar o desempenho com os anos anteriores.
\end{multicols}

\figuramarcadora
\legenda{Legenda em Outfit 9 pt, uma linha abaixo da figura. \fonte{Fonte: LINA, 2026.}}

\subsection{1.1 · Indicadores do ano}
\begin{tabela}{YNN}
  Indicador & \destaque{2025} & \destaque{2026} \\
  Artigos em periódicos & \numero{—} & \numero{—} \\
  Orientações concluídas & \numero{—} & \numero{—} \\
\end{tabela}

\paginademo

% ---------- Finalização ----------
\section{5 · Considerações finais}
\begin{multicols}{2}
Feche o documento com uma síntese do ano em dois ou três parágrafos: o que foi alcançado, o que ficou pendente e as prioridades para o próximo ciclo. Não introduza dados novos nesta seção.

Os anexos, quando houver, vêm antes desta página e seguem a mesma grade das páginas internas, numerados como A, B, C.
\end{multicols}

\subsection{Equipe}
\begin{metadados}{YYY}
  \rotulometa{Coordenação}\newline Nome Sobrenome &
  \rotulometa{Vice-coordenação}\newline Nome Sobrenome &
  \rotulometa{Pesquisadores}\newline Nome Sobrenome\newline Nome Sobrenome \\
  \rotulometa{Pós-graduação}\newline Nome Sobrenome\newline Nome Sobrenome &
  \rotulometa{Graduação}\newline Nome Sobrenome\newline Nome Sobrenome &
  \rotulometa{Técnico}\newline Nome Sobrenome
\end{metadados}

\vfill
\colofao{\destaque{Laboratório de Modelagem da Informação}\newline Departamento de Arquitetura, Urbanismo e Design\newline Universidade Federal do Ceará\newline Av. da Universidade, Fortaleza – CE}
  {contato@lina.ufc.br\newline lina.ufc.br}

\end{document}
